\documentclass[11pt, a4paper]{article}

% Gemeinsame Style-Definitionen (TEXINPUTS muss templates/ enthalten — siehe scripts/build_tex.py)
% bfw-style.tex — gemeinsame Style-Definitionen für alle Handouts/Aufgabenhefte
% Voraussetzung: Kompilation mit xelatex (wegen fontspec)

\usepackage[ngerman]{babel}
\usepackage{geometry}
\geometry{a4paper, top=2.2cm, bottom=2.2cm, left=2.3cm, right=2.3cm}

\usepackage{fontspec}
% Fallback-Kette: Source Sans Pro -> Calibri -> Segoe UI -> Latin Modern Sans
% (Latin Modern ist immer mit MiKTeX vorhanden)
\IfFontExistsTF{Source Sans Pro}{%
  \setmainfont{Source Sans Pro}[Ligatures=TeX]%
}{\IfFontExistsTF{Calibri}{%
  \setmainfont{Calibri}[Ligatures=TeX]%
}{\IfFontExistsTF{Segoe UI}{%
  \setmainfont{Segoe UI}[Ligatures=TeX]%
}{%
  \setmainfont{Latin Modern Sans}[Ligatures=TeX]%
}}}
\IfFontExistsTF{Source Code Pro}{%
  \setmonofont{Source Code Pro}[Scale=0.92]%
}{\IfFontExistsTF{Consolas}{%
  \setmonofont{Consolas}[Scale=0.92]%
}{%
  \setmonofont{Latin Modern Mono}[Scale=0.92]%
}}

\usepackage{xcolor}
% Farbpalette: hell, freundlich, modern
\definecolor{bfwPrimary}{HTML}{0EA5A4}    % Petrol/Türkis - Primär
\definecolor{bfwPrimaryDark}{HTML}{0F766E} % dunkler Petrol für Headings
\definecolor{bfwAccent}{HTML}{F59E0B}     % Amber - Akzent
\definecolor{bfwSuccess}{HTML}{10B981}    % Grün - Erfolg / Lösung
\definecolor{bfwWarn}{HTML}{F97316}       % Orange - Achtung
\definecolor{bfwInk}{HTML}{0F172A}        % fast Schwarz
\definecolor{bfwInkSoft}{HTML}{475569}    % gedämpftes Grau
\definecolor{bfwBgSoft}{HTML}{F8FAFC}     % off-white
\definecolor{bfwBgTeal}{HTML}{ECFEFF}     % helles Türkis
\definecolor{bfwBgAmber}{HTML}{FEF3C7}    % helles Gelb
\definecolor{bfwBgRose}{HTML}{FFE4E6}     % helles Rosé für Eselsbrücken

\usepackage[colorlinks=true, linkcolor=bfwPrimaryDark, urlcolor=bfwPrimaryDark]{hyperref}
\usepackage{enumitem}
\setlist[itemize]{leftmargin=*, itemsep=2pt, topsep=2pt}
\setlist[enumerate]{leftmargin=*, itemsep=2pt, topsep=2pt}

\usepackage[most]{tcolorbox}
\usepackage{booktabs}
\usepackage{tikz}
\usetikzlibrary{decorations.pathmorphing, positioning, shapes.geometric, arrows.meta}

% Merkkästchen — wichtige Definition / Kernpunkt
\newtcolorbox{merksatz}[1][]{%
  enhanced, breakable,
  colback=bfwBgTeal,
  colframe=bfwPrimary,
  boxrule=0.6pt,
  arc=4pt,
  left=10pt, right=10pt, top=8pt, bottom=8pt,
  fonttitle=\bfseries\color{bfwPrimaryDark},
  title={\faLightbulb\ Merksatz},
  #1
}

% Eselsbrücke — Mnemoniks
\newtcolorbox{eselsbruecke}[1][]{%
  enhanced, breakable,
  colback=bfwBgRose,
  colframe=bfwAccent,
  boxrule=0.6pt,
  arc=4pt,
  left=10pt, right=10pt, top=8pt, bottom=8pt,
  fonttitle=\bfseries\color{bfwWarn},
  title={Eselsbrücke},
  #1
}

% Beispiel-Box
\newtcolorbox{beispiel}[1][]{%
  enhanced, breakable,
  colback=bfwBgSoft,
  colframe=bfwInkSoft,
  boxrule=0.4pt,
  arc=3pt,
  left=10pt, right=10pt, top=8pt, bottom=8pt,
  fonttitle=\bfseries\color{bfwInk},
  title={Beispiel},
  #1
}

% Lösungs-Box (für Aufgabenhefte)
\newtcolorbox{loesung}[1][]{%
  enhanced, breakable,
  colback=bfwBgAmber!40,
  colframe=bfwSuccess,
  boxrule=0.6pt,
  arc=3pt,
  left=10pt, right=10pt, top=8pt, bottom=8pt,
  fonttitle=\bfseries\color{bfwSuccess!70!black},
  title={Lösung},
  #1
}

% Glossar-Eintrag
\newcommand{\glossareintrag}[2]{%
  \par\noindent\textbf{\color{bfwPrimaryDark}#1} \enspace #2\par\smallskip
}

% Section-Stil — etwas farbiger als Standard
\usepackage{titlesec}
\titleformat{\section}
  {\Large\bfseries\color{bfwPrimaryDark}}
  {\thesection}{0.6em}{}
\titleformat{\subsection}
  {\large\bfseries\color{bfwPrimary}}
  {\thesubsection}{0.5em}{}
\titleformat{\subsubsection}
  {\normalsize\bfseries\color{bfwInk}}
  {\thesubsubsection}{0.4em}{}

% TikZ-Stil "skizzenhaft" — etwas weicher als Rough.js, aber stabil
\tikzset{
  roughbox/.style = {
    draw=bfwPrimaryDark, thick, fill=bfwBgTeal,
    rounded corners=4pt, inner sep=6pt
  },
  roughaccent/.style = {
    draw=bfwAccent, thick, fill=bfwBgAmber,
    rounded corners=4pt, inner sep=6pt
  },
  roughline/.style = {
    draw=bfwInkSoft, thick, -{Stealth[length=2.5mm]}
  }
}

% Bullet-Symbole (bewusst kein FontAwesome — vermeidet Font-Installations-Probleme)
\newcommand{\faLightbulb}{$\bullet$}
\newcommand{\faBookmark}{$\bullet$}

% Header/Footer
\usepackage{fancyhdr}
\pagestyle{fancy}
\fancyhf{}
\renewcommand{\headrulewidth}{0pt}
\fancyfoot[L]{\small\color{bfwInkSoft}\thepage}
\fancyfoot[R]{\small\color{bfwInkSoft} BFW M-064 Beschaffung im Büromanagement}


\title{\vspace*{-1.5cm}\Huge\bfseries\color{bfwPrimaryDark}Tag~4: Arbeitsschutz \& Gesundheit am Arbeitsplatz\\[0.4em]
  \large\color{bfwInkSoft}\textnormal{Belastung erkennen, Gesundheit schützen}}
\author{\textsf{Büroprozesse gestalten und Arbeitsvorgänge organisieren \textperiodcentered{} M-067}\\
\textsf{\small\copyright\ Can Siebert\,\textperiodcentered\,2026}}
\date{02.07.2026}

\begin{document}
\maketitle
\thispagestyle{empty}

\vspace{1em}
\begin{merksatz}[title={\bfseries Worum geht es heute?}]
Arbeit kann krank machen --- oder gesund halten. Damit Beschäftigte sicher und gesund
arbeiten, verpflichtet der Gesetzgeber die Arbeitgeber zu zahlreichen Schutzmaßnahmen.
In diesem Handout lernen Sie die rechtlichen Grundlagen des Arbeitsschutzes
(ArbSchG, ArbZG, ASiG) und ihre betrieblichen Akteure kennen. Anschließend geht es um
die \emph{Gesundheit am Arbeitsplatz}: Wie wirken Belastungen, wie entstehen Stress,
Burnout und Mobbing --- und mit welchen Maßnahmen lassen sie sich verringern? Das ist
für die Büroarbeit und für die IHK-Abschlussprüfung gleichermaßen wichtig.
\end{merksatz}

\tableofcontents
\vspace{1em}
\hrule
\vspace{1em}

\section{Arbeitsschutz --- der rechtliche Rahmen}

Arbeitgeber sind an verschiedene Gesetze gebunden, die die psychische und physische
Gesundheit der Arbeitnehmer schützen sollen. Dazu gehören das \emph{Arbeitsschutzgesetz},
das \emph{Arbeitssicherheitsgesetz}, das \emph{Arbeitszeitgesetz}, die Gewerbeordnung und
die Arbeitsstättenverordnung. Manche Vorschriften gelten nur für bestimmte Branchen
(z.\,B.\ Bergbau); für Büroarbeitsplätze sind vor allem die allgemeinen Rechtsvorschriften
relevant.

Die Verantwortung für die Erhaltung der Gesundheit am Arbeitsplatz liegt \emph{sowohl beim
Arbeitgeber als auch beim Arbeitnehmer}. Der Arbeitgeber hat dafür zu sorgen, dass
entsprechende Arbeitsbedingungen vorliegen, dass Arbeitsschutzvorschriften und
Unfallverhütungsmaßnahmen eingehalten werden und dass geeignete Arbeitsmittel vorhanden
sind. Der Arbeitnehmer hat ebenfalls die Pflicht, für seine Gesundheit zu sorgen: Er muss
sich so verhalten, dass er Gesundheitsgefährdungen vorbeugt und Gefahren meidet, und den
Arbeitgeber auf Schwachstellen hinweisen.

\begin{merksatz}
Arbeitsschutz ist keine Einbahnstraße: Der Arbeitgeber schafft sichere Bedingungen ---
der Arbeitnehmer muss sie verantwortungsbewusst nutzen.
\end{merksatz}

\subsection{Aufsicht: Berufsgenossenschaft und Gewerbeaufsichtsamt}

Der Arbeitsschutz unterliegt dem Gewerbeaufsichtsamt und den Berufsgenossenschaften.
Jedes Unternehmen gehört branchenspezifisch zu einer \textbf{Berufsgenossenschaft}. Sie
überwacht die Durchführung der Maßnahmen zur Verhütung von Arbeitsunfällen,
Berufskrankheiten und arbeitsbedingten Gesundheitsgefahren, kontrolliert die Erste Hilfe
und berät Unternehmen wie Beschäftigte. Außerdem erlässt sie die \emph{Berufsgenossen\-schaftlichen
Vorschriften (BGV)} zur Unfallverhütung (früher: Unfallverhütungsvorschriften).

Das zuständige \textbf{Gewerbeaufsichtsamt} kontrolliert, ob alle Vorschriften zum
Arbeits-, Umwelt- und Verbraucherschutz eingehalten werden. Es darf Arbeitsstätten
jederzeit betreten und besichtigen. Bei Verstößen kann es Anordnungen erteilen, Betriebe
oder Betriebsteile stilllegen und Bußgelder verhängen (\S~35 GewO).

\subsection{Arbeitszeitgesetz (ArbZG)}

Das \textbf{Arbeitszeitgesetz} soll durch seine Vorschriften den Gesundheitsschutz der
Arbeitnehmer sichern. Festgelegt werden darin beispielsweise die tägliche
\emph{Höchstarbeitszeit} und die \emph{Mindestruhepausen}. Besonderer Schutz gilt für
Nachtarbeiter und Nachtarbeiterinnen. Grundsätzlich besteht ein Arbeitsverbot für Sonn-
und Feiertage, wobei unter besonderen Voraussetzungen Ausnahmen möglich sind.

\subsection{Arbeitsschutzgesetz (ArbSchG)}

Das \textbf{Arbeitsschutzgesetz (ArbSchG)} befasst sich mit der Sicherheit und dem
Gesundheitsschutz der Beschäftigten. Am Anfang des Gesetzes stehen allgemeine Grundsätze
(\S~4 ArbSchG), die für alle Unternehmen gelten:

\begin{enumerate}
  \item Die Arbeit ist so zu gestalten, dass eine Gefährdung für Leben sowie physische und psychische Gesundheit möglichst vermieden und gering gehalten wird.
  \item Gefahren sind an ihrer Quelle zu bekämpfen.
  \item Stand von Technik, Arbeitsmedizin und Hygiene sind zu berücksichtigen.
  \item Technik, Arbeitsorganisation, soziale Beziehungen und Umwelteinfluss sind sachgerecht zu verknüpfen.
  \item Individuelle Schutzmaßnahmen sind nachrangig gegenüber anderen Maßnahmen.
  \item Besonders schutzbedürftige Beschäftigtengruppen sind zu berücksichtigen.
\end{enumerate}

\subsubsection{Gefährdungsbeurteilung (\S 5 ArbSchG)}

Damit der Arbeitgeber nach diesen Grundsätzen handeln kann, muss er nach \S~5 ArbSchG die
Arbeitsbedingungen analysieren und eine \emph{Gefährdungsbeurteilung} erstellen. Dabei
werden ausdrücklich auch die \emph{psychischen Belastungen} berücksichtigt. Eine
Gefährdung kann sich u.\,a.\ ergeben durch die Gestaltung der Arbeitsstätte, physikalische,
chemische und biologische Einwirkungen, Arbeitsmittel, Arbeits- und Fertigungsverfahren
sowie die Arbeitszeit, unzureichende Qualifikation und Unterweisung sowie psychische
Belastung bei der Arbeit.

Die Gefährdungsbeurteilung muss durchgeführt und \emph{dokumentiert} werden:

\begin{itemize}
  \item vor Aufnahme der Tätigkeiten (Erstbeurteilung an allen bestehenden Arbeitsplätzen),
  \item beim Einrichten und Betreiben von Arbeitsstätten,
  \item vor der erstmaligen Verwendung eines Arbeitsmittels.
\end{itemize}

Überprüft und aktualisiert werden muss sie z.\,B., wenn neue Arbeitsplätze eingerichtet
oder neue Arbeitsverfahren eingeführt werden.

\subsubsection{Unterweisung (\S 12 ArbSchG) und Betriebsanweisung}

Damit sich die Mitarbeiter sicherheitsgerecht verhalten können, müssen Unternehmer
regelmäßig \emph{Unterweisungen} während der Arbeitszeit vornehmen. Inhalte sind alle bei
einer Tätigkeit auftretenden Gefahren sowie das richtige Verhalten; allgemeine Themen sind
Brandschutz, Erste Hilfe, Notausgänge und Sammelpunkte. Die Unterweisung erfolgt mündlich,
arbeitsplatzbezogen und praxisnah --- so, dass auch ausländische Arbeitnehmer sie verstehen.
Inhalt und Zeitpunkt müssen schriftlich festgehalten und von der unterwiesenen Person mit
Unterschrift bestätigt werden.

Die \emph{Erstunterweisung} findet vor Aufnahme einer Tätigkeit statt (z.\,B.\ bei
Neueinstellung, Versetzung, neuer Tätigkeit, Arbeit mit Gefahrstoffen). Die Wiederholung
erfolgt mindestens einmal jährlich, bei aktuellem Anlass (z.\,B.\ nach einem Beinah-Unfall)
auch häufiger.

\begin{beispiel}
Martin Küster ist neu in der Produktion eingestellt. In seiner Unterweisung wird er u.\,a.\
in die Arbeit mit der Kreissäge und deren Schutzvorkehrungen eingewiesen. Außerdem werden
ihm Notausgang und Sammelpunkt gezeigt.
\end{beispiel}

Zusätzlich sind \textbf{Betriebsanweisungen} ein wichtiges Instrument. Sie werden vom
Unternehmer oder einem Beauftragten erstellt und am Arbeitsplatz ausgehängt. Darin steht,
welche Gefahren für Mensch und Umwelt bestehen, welche Schutzmaßnahmen zu treffen und
welche Verhaltensregeln einzuhalten sind --- ergänzt um Hinweise für den Gefahrenfall, die
Erste Hilfe und die sachgerechte Entsorgung.

\subsection{Arbeitssicherheitsgesetz (ASiG) und Akteure}

Das „Gesetz über Betriebsärzte, Sicherheitsingenieure und andere Fachkräfte für
Arbeitssicherheit”, kurz \textbf{Arbeitssicherheitsgesetz (ASiG)}, bestimmt, dass der
Arbeitgeber \emph{Betriebsärzte} und \emph{Fachkräfte für Arbeitssicherheit} bestellen muss.
Diese sollen den Arbeitgeber beim Arbeitsschutz und bei der Unfallverhütung unterstützen,
da die Gefahren am wirksamsten aus dem Betrieb heraus bekämpft werden können.

Der \emph{Betriebsarzt} berät den Arbeitgeber (z.\,B.\ bei der Gestaltung von
Arbeitsplätzen, Arbeitsrhythmus und Pausenregelung), untersucht und berät die Arbeitnehmer
arbeitsmedizinisch und beobachtet die Durchführung des Arbeitsschutzes. Zu seinen Aufgaben
gehört es \emph{nicht}, Krankmeldungen auf ihre Berechtigung zu überprüfen.

Zusätzlich schreibt das Sozialgesetzbuch (SGB) vor, dass Unternehmen
\textbf{Sicherheitsbeauftragte} schriftlich bestellen müssen. Sie unterstützen vor Ort als
Bindeglied zwischen Beschäftigten und Arbeitsschutz.

\begin{eselsbruecke}
\textbf{„ArbZG = Zeit, ArbSchG = Schutz, ASiG = Sicherheitspersonal.”} So unterscheiden Sie
die drei wichtigsten Gesetze: Das \emph{Z}eitgesetz regelt die Arbeits\emph{z}eit, das
\emph{Sch}utzgesetz die Schutzmaßnahmen, das ASiG die Bestellung von \emph{Si}cherheits\-fachkräften
und Betriebsärzten.
\end{eselsbruecke}

\section{Gesundheit am Arbeitsplatz}

Die Gesundheit eines Menschen ist ein wichtiges Gut. Neben dem staatlichen
Gesundheitsschutz durch Gesetze bieten viele Unternehmen eigene Angebote (z.\,B.\
Betriebssport, Kooperation mit Fitnessstudios), um die Gesundheit zu erhalten und Stress
abzubauen --- das \emph{betriebliche Gesundheitsmanagement (BGM)}.

\subsection{Belastung und Beanspruchung}

\textbf{Belastungen} bringt jede Tätigkeit mit sich. Der Begriff ist in der
Arbeitswissenschaft zunächst \emph{neutral} und bezeichnet alle Einflüsse, die von außen
auf den Menschen einwirken. Wie stark sich jemand durch eine Belastung \emph{beansprucht}
fühlt, hängt von verschiedenen Faktoren ab: Belastungsstärke, Alter, Konstitution,
Qualifikation, Entscheidungs- und Handlungsspielraum sowie dem persönlichen Umgang mit
Stress. Sowohl Über- als auch Unterforderung führen zu \emph{Fehlbelastungen}.

\begin{merksatz}
\textbf{Belastung} ist der Einfluss von außen (neutral). \textbf{Beanspruchung} ist die
individuelle Wirkung im Menschen. Dieselbe Belastung wirkt je nach Person ganz
unterschiedlich.
\end{merksatz}

Seit Oktober 2013 werden \emph{psychische Belastungen} ausdrücklich als
Gesundheitsgefährdung angesehen und im Arbeitsschutzgesetz berücksichtigt. Jeder
Arbeitgeber muss sie durch die Gefährdungsbeurteilung feststellen und entsprechende
Maßnahmen ergreifen. Wichtige psychische Belastungsfaktoren und Gegenmaßnahmen:

\begin{center}
\begin{tabular}{p{4.2cm} p{5.0cm} p{4.2cm}}
\toprule
\textbf{Bereich} & \textbf{Belastungsfaktor} & \textbf{Gegenmaßnahme}\\
\midrule
Soziale Beziehungen & Fehlende Rückmeldung; Leistung nicht gewürdigt & Feedbackgespräche, Lob\\
Arbeitsaufgabe & Über- oder Unterforderung & Angemessene Personal- und Aufgabenplanung\\
Arbeitsumgebung & Klima, Licht, Lärm & Gesetzliche Bestimmungen anpassen\\
Betriebliche Organisation & Ständige Unterbrechung, unklare Kompetenzen & Prioritäten setzen, Kompetenzen klar regeln\\
\bottomrule
\end{tabular}
\end{center}

\subsection{Stress}

\textbf{Stress} ist eine natürliche Reaktion des Körpers auf eine Herausforderung --- ein
Überbleibsel aus den Anfangszeiten der Menschheit. Bei Gefahr bereitete sich der Körper auf
\emph{Kampf oder Flucht} vor: Puls und Blutdruck steigen, die Sinne werden geschärft, die
Muskeln angespannt, Stresshormone werden ausgeschüttet. Heute können wir am Arbeitsplatz
weder kämpfen noch fliehen --- der innere Druck wird nicht abgebaut. Der Körper bleibt in
ständiger Alarmbereitschaft, und es droht die Gefahr ernsthafter Erkrankungen.

\begin{center}
\begin{tabular}{p{6.5cm} p{6.5cm}}
\toprule
\textbf{Eustress (positiver Stress)} & \textbf{Distress (negativer Stress)}\\
\midrule
Erhöht die Aufmerksamkeit; fördert die Leistungsfähigkeit ohne zu schaden; motiviert; steigert die Produktivität & Entsteht, wenn Stress zu häufig oder ohne körperlichen Ausgleich auftritt; überfordert; ist bedrohlich\\
\bottomrule
\end{tabular}
\end{center}

\medskip
Symptome bei negativem Stress sind z.\,B.\ verminderte Leistungs- und
Konzentrationsfähigkeit, Schlafstörungen, Magen-Darm-Probleme, Gereiztheit und
Bluthochdruck. Ständiger Stress schwächt das Immunsystem und begünstigt Herzerkrankungen
und Rückenbeschwerden; auch Burnout kann eine Folge sein.

\begin{merksatz}
\textbf{Stress abbauen:} Sport treiben, Freizeitstress vermeiden, Erholungsphasen
einplanen \emph{bevor} man sich gestresst fühlt, ausreichend schlafen und Entspannung
lernen (z.\,B.\ Yoga, Meditation). Bewegung beschleunigt den Stressabbau.
\end{merksatz}

\subsection{Burnout}

Der Begriff \textbf{Burnout} kommt aus dem Englischen und bedeutet „ausbrennen”. Gemeint ist
ein Zustand völliger physischer und psychischer Erschöpfung aufgrund meist beruflicher
Überlastung. Mediziner sehen im Burnout eine \emph{Vorstufe zur Depression}. Burnout ist
keine eigenständige Krankheit mit genau definierter Diagnose und entsteht schleichend.

Typische Symptome (über mehrere Wochen): chronische Müdigkeit, Konzentrationsprobleme,
Nervosität, Schlafstörungen, Erschöpfung, Kopf- und Rückenschmerzen, Verspannungen,
Schwindel, Hörsturz. Besonders gefährdet sind Menschen mit hohen Leistungsansprüchen,
Perfektionisten und Personen, die sich vor allem über ihren Job definieren. Ursachen sind
oft großes Engagement und übertriebener Ehrgeiz.

Mithilfe der \textbf{Work-Life-Balance} sollen private und berufliche Anforderungen in ein
gesundes Gleichgewicht gebracht werden. Maßnahmen der Unternehmen sind z.\,B.\ flexible
Arbeitszeitmodelle (Teilzeit, Gleitzeit, Vertrauensarbeitszeit), veränderte
Arbeitsorganisation (Job Rotation, Job Enlargement), betriebliches Gesundheitsmanagement,
Unterstützung bei Alltagsproblemen (Kinder-/Pflegebetreuung) und Weiterbildung im Stress-,
Zeit- und Konfliktmanagement.

\subsection{Mobbing}

Das englische Verb „to mob” bedeutet „lärmend herfallen über, anpöbeln, angreifen,
attackieren”. Nach der Definition von Bernd Zuschlag beschreibt \textbf{Mobbing} ein
schikanöses Handeln einer oder mehrerer Personen, das gegen eine Einzelperson oder
Personengruppe gerichtet ist. Die schikanösen Handlungen werden meist über einen längeren
Zeitraum wiederholt; sie implizieren die Täterabsicht, das Opfer zu schädigen.

Im Unterschied zu gewöhnlichen Konflikten ist Mobbing an weiteren Kriterien zu erkennen:

\begin{itemize}
  \item Der Konflikt hat sich verfestigt.
  \item Die Würde des Menschen wird angegriffen, dem Opfer wird das positive Selbstwertgefühl genommen.
  \item Die Konfliktparteien begegnen sich nicht auf Augenhöhe.
  \item Der Unterlegene kommt nicht aus eigener Kraft aus der Situation heraus.
\end{itemize}

\emph{Kein} Mobbing liegt vor bei einer einmaligen Aktion, einer notwendigen
organisatorischen Maßnahme oder wenn keine schikanöse Absicht dahintersteckt.

\subsubsection{Ursachen und Folgen}

Die Ursachen können beim \emph{Täter} (Ziele wie Machtausübung; Ängste wie Autoritätsverlust),
beim \emph{Opfer} (besonders gefährdet: neue Mitarbeiter, Konkurrenten um eine Position,
befristet Beschäftigte; Anlass berechtigt oder unabsichtlich) oder in den
\emph{Rahmenbedingungen} der Arbeitswelt liegen (Leistungsdruck, umständliche Dienstwege,
Kündigungsregelungen, abweichende soziale Normen).

Die \emph{Folgen} reichen vom Mobbing-Opfer (körperliche und psychische Beschwerden bis hin
zu Selbstmordgedanken, Leistungsabfall, „innere Kündigung”) über den Täter (Stress,
Erkrankungsgefahr) und den Betrieb (schlechtes Betriebsklima, Fehlzeiten, Fluktuation) bis
zur Volkswirtschaft (jährliche Kosten in Milliardenhöhe).

\subsubsection{Maßnahmen gegen Mobbing}

Voraussetzung ist zunächst die Feststellung der Ursachen. Das Opfer kann eine
\emph{selbstkritische Prüfung} vornehmen, den \emph{Täter} auf sein Tun ansprechen und
ihm Grenzen setzen, das \emph{soziale Netzwerk} und das \emph{Selbstbewusstsein} stärken,
ein \emph{Mobbing-Tagebuch} führen, den \emph{Rechtsweg} beschreiten oder als letzte
Möglichkeit den \emph{Arbeitsplatz wechseln}. Entscheidend ist, früh zu reagieren. Auch
Unternehmen können vorbeugen: durch klare Arbeitsorganisation, Förderung der
Arbeitszufriedenheit und eine an ethischen Grundsätzen orientierte Unternehmensphilosophie.

\subsection{Suchtprävention und gesunde Ernährung}

\textbf{Suchtprävention} gehört in das betriebliche Gesundheitsmanagement: Faktoren wie
hoher Leistungsdruck, ungünstige Arbeitszeiten, schwieriges Betriebsklima oder ein
unregelmäßiger Schlafrhythmus können das Suchtrisiko erhöhen. Mitarbeiter sollten
informiert, Führungskräfte geschult und Beratungsgespräche angeboten werden. Wichtig:
Alkoholprobleme von Kollegen dürfen nicht verschwiegen oder gedeckt werden. Suchtmittel
beeinträchtigen die Arbeitssicherheit (Konzentrationsstörungen, längere Reaktionszeit,
getrübtes Urteilsvermögen, Unfallgefahr).

\textbf{Gesunde Ernährung im Büro} hält fit und leistungsfähig --- das Gehirn verbraucht
ca.\ 15\,\% des täglichen Energiebedarfs. Empfehlenswert sind ein ausgewogenes Frühstück
(Vollkornbrot, magerer Käse, Obst), fettarme und langsam verdauliche Zwischenmahlzeiten in
einer bewussten Pause sowie ein kalorienarmes Mittagessen (Salate, Gemüse, Kartoffeln,
Vollkornprodukte). Wichtig ist ausreichendes Trinken (Wasser, ungesüßter Tee, Saftschorle
im Verhältnis 1/3 Saft zu 2/3 Wasser); Kaffee und Alkohol sind Genussmittel und sollten nur
in kleinen Mengen konsumiert werden.

\section{Zusammenfassung}

\begin{enumerate}
  \item \textbf{Arbeitsschutz} schützt Sicherheit und Gesundheit; zentrale Gesetze sind ArbSchG, ArbZG und ASiG.
  \item Verantwortung tragen \textbf{Arbeitgeber und Arbeitnehmer}; Aufsicht führen Berufsgenossenschaft (BGV) und Gewerbeaufsichtsamt (\S~35 GewO).
  \item \textbf{Gefährdungsbeurteilung} (\S~5), \textbf{Unterweisung} (\S~12) und \textbf{Betriebsanweisung} sind die zentralen Instrumente im Betrieb.
  \item \textbf{Belastung} ist neutral, \textbf{Beanspruchung} ist die individuelle Wirkung; psychische Belastungen werden seit 2013 ausdrücklich berücksichtigt.
  \item \textbf{Stress}: Eustress hilft, Distress schadet; Bewegung und Erholung bauen Stress ab. \textbf{Burnout} ist eine Vorstufe zur Depression.
  \item \textbf{Mobbing} ist wiederholte Schikane; früh reagieren. \textbf{BGM} umfasst Suchtprävention, gesunde Ernährung und Work-Life-Balance.
\end{enumerate}

\newpage

\section*{Glossar}
\addcontentsline{toc}{section}{Glossar}

\glossareintrag{Arbeitsschutzgesetz (ArbSchG)}{Gesetz mit Maßnahmen, die Sicherheit und Gesundheitsschutz der Beschäftigten bei der Arbeit sichern und verbessern.}

\glossareintrag{Arbeitszeitgesetz (ArbZG)}{Regelt zum Gesundheitsschutz u.\,a.\ Höchstarbeitszeit, Ruhepausen, Nachtarbeit und das Sonn-/Feiertagsarbeitsverbot.}

\glossareintrag{Arbeitssicherheitsgesetz (ASiG)}{Verpflichtet den Arbeitgeber, Betriebsärzte und Fachkräfte für Arbeitssicherheit zu bestellen.}

\glossareintrag{Gefährdungsbeurteilung}{Beurteilung der Gefährdungen am Arbeitsplatz (\S~5 ArbSchG); ermittelt notwendige Schutzmaßnahmen, dokumentiert, auch psychische Belastung.}

\glossareintrag{Unterweisung}{Regelmäßige, arbeitsplatzbezogene Belehrung der Beschäftigten (\S~12 ArbSchG); schriftlich festgehalten und unterschrieben.}

\glossareintrag{Betriebsanweisung}{Schriftliches Dokument zu Gefahren, Schutzmaßnahmen und Verhaltensregeln; am Arbeitsplatz ausgehängt.}

\glossareintrag{Berufsgenossenschaft}{Branchenspezifischer Unfallversicherungsträger; überwacht Unfallverhütung und erlässt die BGV.}

\glossareintrag{Gewerbeaufsichtsamt}{Staatliche Aufsicht über Arbeits-, Umwelt- und Verbraucherschutz; kann Betriebe stilllegen und Bußgelder verhängen.}

\glossareintrag{Sicherheitsbeauftragte}{Vom Unternehmen schriftlich bestellte Personen (SGB), die den Arbeitsschutz vor Ort unterstützen.}

\glossareintrag{Belastung}{Arbeitswissenschaftlich neutraler Begriff für alle Einflüsse, die von außen auf den Menschen einwirken.}

\glossareintrag{Beanspruchung}{Die individuelle Wirkung einer Belastung im Menschen; abhängig u.\,a.\ von Konstitution und Stressumgang.}

\glossareintrag{Eustress}{Positiver Stress; erhöht Aufmerksamkeit und Leistungsfähigkeit, motiviert.}

\glossareintrag{Distress}{Negativer Stress; entsteht bei zu häufigem Stress ohne Ausgleich, überfordert und ist bedrohlich.}

\glossareintrag{Burnout}{Zustand völliger körperlicher und psychischer Erschöpfung durch berufliche Überlastung; Vorstufe zur Depression.}

\glossareintrag{Work-Life-Balance}{Gesundes Gleichgewicht zwischen beruflichen und privaten Anforderungen.}

\glossareintrag{Mobbing}{Über längere Zeit wiederholtes schikanöses Handeln gegen eine Person mit Schädigungsabsicht.}

\glossareintrag{Betriebliches Gesundheitsmanagement (BGM)}{Betriebliche Maßnahmen zur Erhaltung der Gesundheit, z.\,B.\ Betriebssport, Suchtprävention, Ernährung.}

\glossareintrag{Suchtprävention}{Maßnahmen im BGM zur Vorbeugung von Suchterkrankungen (Information, Schulung, Beratung).}

\end{document}
